\begin{proposition}[The moving-gap Green kernel]
\label{prop:gap-kernel-green}
Let
\[
 P(t)=(2t+1)(17t^2+17t+5),
\]
and let $U_m(a)$ be the restart continuants defined by
\[
 U_{-1}(a)=0,\qquad U_0(a)=1,\qquad
 U_{m+1}(a)=P(a+m)U_m(a)-(a+m)^6U_{m-1}(a).
\]
For an indeterminate~$a$, put
\[
 V_m(a)=\frac{U_m(a)}{(a+1)_m^3},\qquad
 (a+1)_m=\prod_{j=1}^m(a+j),\qquad V_{-1}(a)=0,
\]
and set
\[
 F_a(z)=\sum_{m\geq0}V_m(a)z^m,
 \qquad \theta=z\frac{d}{dz}.
\]
Then the following assertions hold in $\mathbb Q(a)[[z]]$.

\begin{enumerate}
\item The normalized continuants satisfy
\begin{equation}
 (a+m+1)^3V_{m+1}(a)
 =P(a+m)V_m(a)-(a+m)^3V_{m-1}(a)
 \qquad(m\geq0).
 \label{eq:green-normalized-recurrence}
\end{equation}

\item If
\begin{equation}
 \mathscr D_a
 = (\theta+a)^3-zP(\theta+a)
   +z^2(\theta+a+1)^3,
 \label{eq:green-operator}
\end{equation}
then
\begin{equation}
 \boxed{\mathscr D_aF_a=a^3.}
 \label{eq:green-inhomogeneous}
\end{equation}
Equivalently, if
\[
 \mathscr L_0=\theta^3-zP(\theta)+z^2(\theta+1)^3,
\]
then, formally under the conjugation
$\theta z^a=z^a(\theta+a)$,
\begin{equation}
 \boxed{\mathscr L_0\bigl(z^aF_a(z)\bigr)=a^3z^a.}
 \label{eq:green-conjugated}
\end{equation}
Thus $z^aF_a$ is a Green-kernel column with monomial source, not an
eigenfunction of~$\mathscr L_0$.

\item The coefficient of $\theta^3$ in~$\mathscr D_a$ is
\[
 Q(z)=1-34z+z^2.
\]
Consequently~\eqref{eq:green-inhomogeneous}, after adjoining the constant
solution~$1$, is a homogeneous first-order system of rank four whose pole
divisor on the $z$-line is contained in
\begin{equation}
 z=0,\qquad z^2-34z+1=0,\qquad z=\infty,
 \label{eq:green-singular-divisor}
\end{equation}
independently of~$a$.  Its coefficients have bounded degree in~$a$.

\item Let $b,c$ be the two solutions of the Ap\'ery recurrence with
$(b_0,b_1)=(1,5)$ and $(c_0,c_1)=(0,1)$.  For every integer $a\geq1$ and
$m\geq0$,
\begin{equation}
 \boxed{
 V_m(a)=a^3\bigl(b_{a-1}c_{a+m}-c_{a-1}b_{a+m}\bigr).}
 \label{eq:green-casoratian}
\end{equation}
In particular,
\[
 V_m(0)=b_m,\qquad V_m(1)=c_{m+1}.
\]
\end{enumerate}
\end{proposition}

\begin{proof}
Dividing the recurrence for $U_{m+1}(a)$ by
$(a+1)_{m+1}^3$ gives~\eqref{eq:green-normalized-recurrence}; the second
term uses
\[
 \frac{(a+m)^6(a+1)_{m-1}^3}{(a+1)_{m+1}^3}
 =\frac{(a+m)^3}{(a+m+1)^3}.
\]

Multiply~\eqref{eq:green-normalized-recurrence} by $z^{m+1}$ and sum over
$m\geq0$.  Since $V_0=1$ and $V_{-1}=0$, the three resulting sums are
\begin{align*}
 \sum_{m\geq0}(a+m+1)^3V_{m+1}z^{m+1}
  &=(\theta+a)^3F_a-a^3,\\
 \sum_{m\geq0}P(a+m)V_mz^{m+1}
  &=zP(\theta+a)F_a,\\
 \sum_{m\geq0}(a+m)^3V_{m-1}z^{m+1}
  &=z^2(\theta+a+1)^3F_a.
\end{align*}
This proves~\eqref{eq:green-inhomogeneous}.  Conjugating by $z^a$ gives
\eqref{eq:green-conjugated}.

The leading coefficient in~$\theta$ is
$1-34z+z^2$, because the leading coefficient of $P$ is~$34$.
Writing
\[
 Y=(1,F_a,\theta F_a,\theta^2F_a)^{\mathsf T}
\]
and solving~\eqref{eq:green-inhomogeneous} for $\theta^3F_a$ gives a
first-order equation $\theta Y=A(a,z)Y$ with entries in
$\mathbb Q(a,z)$ and denominators dividing~$Q(z)$.  Passing from
$\theta$ to $d/dz$ introduces only the additional factor~$z$; this proves
the pole-divisor assertion.  The bounded-degree assertion is immediate
from~\eqref{eq:green-operator}.

For the final identity, define
\[
 K_m=a^3\bigl(b_{a-1}c_{a+m}-c_{a-1}b_{a+m}\bigr).
\]
The Ap\'ery recurrence at index $a+m$ shows that $K_m$ satisfies
\eqref{eq:green-normalized-recurrence}.  The companion Casoratian gives
\[
 b_{a-1}c_a-c_{a-1}b_a=\frac1{a^3},
\]
so $K_0=1$; the recurrence at index~$a$, together with
$K_{-1}=0$, gives $K_1=P(a)/(a+1)^3$.  These are the initial values of
$V_m(a)$, proving~\eqref{eq:green-casoratian}.  Finally,
$U_m(0)=(m!)^3b_m$ gives $V_m(0)=b_m$, while
\eqref{eq:green-casoratian} at $a=1$ gives $V_m(1)=c_{m+1}$.
\end{proof}

\begin{corollary}[Long gaps as coefficients of the Green kernel]
\label{cor:long-gap-green}
For integers $G\geq1$ and $u\geq0$,
\begin{equation}
 N_G(u)=U_{G-1}(u+1)
 =\bigl((u+2)_{G-1}\bigr)^3V_{G-1}(u+1).
 \label{eq:long-gap-green}
\end{equation}
If $0\leq u<u+G<p$, all factors in the rising factorial are units
modulo~$p$, and hence
\[
 N_G(u)=0\pmod p
 \quad\Longleftrightarrow\quad
 [z^{G-1}]F_{u+1}(z)=0\pmod p.
\]
Equivalently,
\[
 \frac{N_G(u)}{\prod_{j=1}^G(u+j)^3}
 =b_uc_{u+G}-c_ub_{u+G}.
\]
\end{corollary}

\begin{proof}
The first identity is the restart dictionary and the definition of~$V$.
The nonwrapping assumption makes every factor $u+2,\ldots,u+G$ nonzero
modulo~$p$, proving the equivalence.  The final identity follows from
\eqref{eq:green-casoratian} with $a=u+1$ and $m=G-1$ after multiplying by
the missing factor $(u+1)^3$.
\end{proof}

\begin{remark}[What the differential compression does not supply]
\label{rem:green-jet-barrier}
The operator~\eqref{eq:green-operator} has fixed order and fixed singular
support, but the bridge length~$G$ remains the coefficient index
$[z^{G-1}]$.  Algebraizing this coefficient by a naive jet construction
has order growing with~$G$.  Thus Proposition~\ref{prop:gap-kernel-green}
does not by itself give a bounded-conductor finite-field family or an
estimate for the long-bridge incidence sum.  It instead isolates the
remaining geometric task: realize the two-index coefficients of this
inhomogeneous Green kernel by a fixed-dimensional arithmetic family, or
prove the required prime-averaged coefficient nonconcentration directly.
\end{remark}
